\begin{trapbox}
\begin{center}
\begin{tikzpicture}[x=.88cm,y=.88cm,font=\small,
  every node/.style={inner sep=3pt},line cap=round]

\foreach \s/\p/\h/\title in {0/4/1/{外点 $P(4,0)$：$S=4$},6.6/1/4/{内点 $P(1,0)$：$S=\frac14$}}{
\begin{scope}[shift={(\s,0)}]
\draw[->,gray!65] (-2.4,0)--(4.6,0);
\draw[CStatement,very thick] (0,0) ellipse [x radius=2,y radius=1];
\draw[CTrap,very thick] (\h,-1.45)--(\h,1.45);
\fill (\p,0) circle (1.6pt) node[below] {$P$};
\node at (1,2) {\title};
\node[CTrap] at (\h,-1.9) {$x=\h$};
\end{scope}}
\draw[orange!85!black,thick] (1,.866025)--(4,0)--(1,-.866025);
\fill (1,.866025) circle (1.6pt); \fill (1,-.866025) circle (1.6pt);

\end{tikzpicture}

{\small 同尺度比较 $\frac{x^2}{4}+y^2=1$：外点有切点弦，内点的形式代换线不交椭圆。}
\end{center}


    \begin{description}[style=nextline, leftmargin=2.8em, labelsep=0.5em, itemsep=0.5em]
        \item[\textbf{\textcolor{CTrap}{易错1：不先判断点的位置}}]
        \textbf{错：} 见到 $P(x_0,y_0)$ 就将乘积式叫作``切点弦''。
        \par\textbf{对：} 先算 $S=\frac{x_0^2}{a^2}+\frac{y_0^2}{b^2}$。
        $S>1$ 才有两个不同切点；$S=1$ 是点 $P$ 处的切线；$S<1$ 没有从 $P$ 引出的实切线。
        若 $P=O$，乘积式是 $0=1$，甚至不是一条直线。

        \item[\textbf{\textcolor{CTrap}{易错2：把切点弦误当成切线}}]
        \textbf{错：} 认为切点弦一定经过 $P$。
        \par\textbf{对：} $PA,PB$ 是两条切线；$AB$ 是两个切点的连线。
        $P$ 代入 $AB$ 方程，左端等于 $S>1$，因此 $P\notin AB$。

        \item[\textbf{\textcolor{CTrap}{易错3：分母忘记平方或符号写错}}]
        \textbf{错：} 写成 $\frac{x_0x}{a}+\frac{y_0y}{b}=1$，或把椭圆第二项写成减号。
        \par\textbf{对：} 对照原椭圆
        $\frac{x^2}{a^2}+\frac{y^2}{b^2}=1$，只替换分子中的平方，分母与加号不变。

        \item[\textbf{\textcolor{CTrap}{易错4：反求外点时直接比较未归一的系数}}]
        \textbf{错：} 已知切点弦 $4x+3y=6$，就直接令
        $x_0/a^2=4,\ y_0/b^2=3$。
        \par\textbf{对：} 先两端除以 $6$：$\frac23x+\frac12y=1$，再比较系数。
        特别地，$Ax+By=0$ 经过椭圆中心，不能是某个有限外点的切点弦。

        \item[\textbf{\textcolor{CTrap}{易错5：证明中出现斜率分母为零}}]
        \textbf{错：} 先把任意切线写成 $y=kx+m$，忽略竖直的直线。
        \par\textbf{对：} 使用方法三的直线表达
        $x=x_1+t(x_0-x_1)$、$y=y_1+t(y_0-y_1)$，可同时处理水平、竖直和倾斜直线。

        \item[\textbf{\textcolor{CTrap}{易错6：把圆与椭圆的几何性质完全等同}}]
        \textbf{错：} 因为圆中 $OM\perp AB$，就断言椭圆中也必有 $OM\perp AB$。
        \par\textbf{对：} 不同方向的缩放保持直线和中点，但不一定保持角度。椭圆情形一般只有 $O,M,P$ 三点共线；不能直接断言垂直。
    \end{description}
\end{trapbox}
